\section{模型评价与推广}

\subsection{模型优点}

\textbf{第一，四个问题建立在同一套方程与同一个数值框架上，结论之间可以互相校验。} 控制方程\eqref{eq:heat}、\eqref{eq:mass} 与边界条件\eqref{eq:robin} 在四问中形式一致，只有物性参数集、求解时段、判据与坐标是否随材料运动有所不同；因此问题二的终点时刻与问题三完全相同（同一模型、仅输出步长不同），问题四与对照算例的差异可以定量归因于几何变化。这种``统一框架 $+$ 局部特化''的结构，使四问的结果具有可比性，也使误差来源可以被分离。

\textbf{第二，问题四的移动边界处理在数学上干净、在实现上简单。} 通过材料坐标变换，把随时间收缩的求解域映射到固定域，并\textbf{证明}了坐标变换引入的输运项与材料速度带来的对流项严格相消（式\eqref{eq:cancel}）。由此得到三个实质好处：方程中不含对流项，故没有 CFL 限制、也不需要迎风格式；不需要动的网格或重映射；半径的分段线性插值在节点处的导数跳变不进入方程，因而不会造成数值困难。这一处理还使质量审计（式\eqref{eq:massaudit}）成为一个\textbf{强有力的正确性判据}：若相消被误用或干基密度幂次写错，残差会立刻从 $10^{-15}$ 跳到 $O(1)$。

\textbf{第三，边界条件在通量层面实现，离散格式天然守恒。} 采用节点中心有限体积、以面通量差表示散度，使内部通量在总体收支中成对抵消，实测质量收支相对残差在 $10^{-15}$ 量级。相比``先写节点方程再补边界''的做法，这一实现方式既避免了轴心 $1/\rpos$ 奇性（轴心通量直接取零），也使符号方向可以用加热与干燥两个极限情形逐一核验。

\textbf{第四，结果的可信度有独立的证据链。} 本文同时给出了四类互相独立的检验：网格加密与 Richardson 估计（空间）、BDF 与 Radau IIA 交叉验证与容差收紧（时间）、圆柱 Bessel 级数解析解对照（物理模型）、离散守恒审计与独立目录复算（实现与格式）。其中解析解对照是最强的一项——它不依赖求解器的任何内部假设。此外，全部交付文件的结构由程序硬校验后才写出，避免了``格式错了但看起来跑通了''的静默失效。

\textbf{第五，结果的敏感性与局限性做了量化。} 环境外推假设的影响被量化为小时量级并与实测波动幅度绑定；网格与时间误差被量化为与四位小数分辨率的关系；每一个结构性假设都写明了代价。这使结论既可以复算，也可以被质疑。

\subsection{模型缺点与不确定性}

\textbf{最主要的缺点是表面热边界不含蒸发潜热。} 由于题目未给出汽化潜热及其与蒸发速率的关系，本文没有引入这一参数，代价是干燥后期的表面温度可能偏高、干燥速率可能偏快。考虑到本题的干燥时长长达 $2\sim3$ 天，这一缺失的影响不容忽视，且无法用题目数据量化。这是模型在\textbf{物理完整性}上最明显的短板。

\textbf{第二项不确定性来自传质边界的基准。} 题给形式 $\Ddiff{\C}_\rpos|_\R=-\hmass({\C}_s-\Cinf)$ 在数值上被直接使用，但题目未说明气相水分的 kg/kg 与药材干基含水率是否同一基准。若两者相差一个数量级，边界驱动势将整数量级地偏，进而使烘干时长整体偏移。本文把它作为``有效闭合假设''如实披露，而不是当作已证明的物理关系。

\textbf{第三项不确定性来自环境数据的外推。} 附件 1 只覆盖前 $4$~h，而问题 2/3/4 要跑到 $2\sim3$ 天，$93\%$ 的时长依赖末值保持假设。灵敏度分析（表~\ref{tab:sensitivity}）给出的量级是 $0.65\sim0.75$~h，小于结构性闭合的影响，但在评审视角下它是最容易被质疑的一环。

\textbf{第四，几何简化带来两处误差。} 一维轴对称忽略了端面效应（端面与侧面积之比 $0.08$），题目未给端面边界条件，该误差无法约束；问题四假设径向均匀收缩且长度不变，而附件 2 只给出半径，无法排除长度同步缩短的可能性（两种情形下体积分别剩 $21.5\%$ 与 $35.9\%$）。

\textbf{第五，问题四的定性结论依赖于闭合形式。} ``收缩使干燥加快''是本模型在干基通量闭合下的推论：干基表观密度按 ${\R}^{-2}$ 增大，抵消并超过了表面积按 $\R$ 减小的影响。若换用按表面积定义的蒸气通量闭合，结论的方向可能改变。本文不声称该结论对闭合形式不敏感。

\textbf{第六，未引入平衡含水率。} 本文把 $\Cinf$ 只当作边界驱动势，因此不排除``含水率可以无限降低''这一模型行为。真实药材存在吸附平衡含水率，若要求低于平衡值的含水率，模型预测的烘干时间会被低估。

\subsection{改进方向}

\textbf{其一，补齐表面能量收支。} 在式\eqref{eq:robin} 的表面热边界中加入蒸发耗热项，即

\begin{equation}
\kcond\left.\frac{\partial \T}{\partial \rpos}\right|_{\R}=\hconv(\Tinf-{\T}_s)-\hmass\,({\C}_s-\Cinf)\,\Delta H_{\mathrm{vap}},
\end{equation}
其中 $\Delta H_{\mathrm{vap}}$ 需查表并注明出处（题目未给出）。这会同时降低表面温度与蒸发速率，方向上是``使干燥变慢''，可作为当前结果的保守性界。

\textbf{其二，把有效传质边界升级为两相模型。} 引入气相水蒸气分压与药材含水率之间的吸附等温线（如 Guggenheim--Anderson--de Boer 模型），把边界条件写成蒸气分压差驱动的形式，并据此重新标定 $\hmass$。这需要补充相对湿度数据，但可以从根本上解决基准不一致的问题。

\textbf{其三，把收缩与干燥耦合起来。} 本模型直接把附件 2 的 $\R(\t)$ 作为输入，而更完整的做法是由失水量反推收缩量（例如假设体积减少等于失去的水体积），使几何变化成为解的一部分。本文做过量级检查：以本模型的解为例，水的体积分数由初值约 $0.71$ 降到结束时刻约 $0.08$，减少 $0.63$；而固定长度假设下材料体积分数由 $1.00$ 降到 $0.36$，减少 $0.64$，两者相差约 $1$ 个百分点，说明附件 2 的 $\R(\t)$ 与失水量在量级上自洽。但这一检查不能当作验证，因为它未验证体积可加性、且用的是本模型自己的解。

\textbf{其四，二维轴对称模型。} 若给出端面边界条件与轴向装料方式，可以把模型扩展到二维轴对称，定量评估端面效应对烘干时长的影响；本文的一维结果是其零阶近似。

\textbf{其五，把模型用于工艺优化。} 在现有框架下，把 $\Tinf(\t)$ 由``实测序列''改为``待优化的控制变量''，以烘干时长、能耗与品质（例如有效成分保留率）为目标构造优化问题，即可回答``如何设计升温曲线''这一工艺问题。由于单次求解只需数十秒（问题三约 \QThreeWallS~s、问题四约 \QFourWallS~s），该优化在计算上是可行的。

\subsection{推广}

本文的模型与算法结构并不依赖于单味药材的具体参数，具有较宽的适用范围。

\textbf{在干燥对象上}，只要物料近似为圆柱形多孔介质、干燥过程由内部扩散控制且含水率可定义为干基质量比，模型即可直接套用；例如中药材中的天麻、桔梗、白术等根茎类药材的热风干燥研究普遍采用圆柱或切片样本与其有效扩散系数模型\cite{gastrodia2021,platycodon2020}，木瓜、秋葵等果蔬的收缩干燥也有类似的处理框架\cite{okra2025}。把附录 2/3/4 换成对应物料的经验公式即可。

\textbf{在干燥方式上}，模型把环境条件抽象为 $\Tinf(\t)$ 与 $\Cinf(\t)$ 两条时间序列，因此同样适用于热风以外的情形：只要给定微波、真空或太阳能干燥的环境温度与湿度时程，甚至分段间歇的工况（这正是``间歇干燥让内部水分重新分布''的工艺思路），都可以用同一套方程求解\cite{review2021,gastrodia2021}。多阶段工况在本文中甚至不需要额外建模——收敛分析表明，把环境条件写成分段函数即可。

\textbf{在问题形式上}，本文把``确定烘干时长''处理成事件检测、把``变物性与移动边界''处理成系数与坐标变换，这两个技巧可以推广到其他``带阈值的枯燥时间''问题：例如药物溶出达到某一释放度的时间、食品冷冻达到中心温度要求的时间、以及材料干燥/固化达到指定含水量的时间。这类问题的共同结构是``一个单调泛函首次穿越阈值''，都可以用一次积分加事件检测解决，而不必做反复的试算。

\textbf{在方法论上}，本文给出的验证链条（守恒审计 $\to$ 网格加密 $\to$ 时间方法交叉验证 $\to$ 解析解对照 $\to$ 独立复算 $\to$ 灵敏度分析）可以作为同类数模问题的工作模板；其中``把结构性局限与数值误差分开陈述''的做法，尤其值得保留：$10^{-15}$ 的守恒残差并不能弥补一个错误的物理闭合。
